\subsection*{A pictured premise helps restore missing Chinese writing}
We trained the same visual writer to assess a human-written claim against a premise and to reconstruct an original phrase removed from that claim. OCNLI provides premises with several claims that express different relations\cite{ocnli}. On 2,947 official development cases the paired writer generated 1,610 complete correct decision sentences, compared with 1,514 when training reassigned premises (difference 3.26 percentage points, 95\% premise-bootstrap interval 1.41--5.08). It completed all three decisions for six of the 35 official premises spanning all three relations. A linear model of visual unigrams, bigrams and their overlaps gave 1,689 correct decisions and ten complete triples. On constructed subject reversals, the paired writer gave identical decisions for 46 of 48 passage pairs and answered both sides correctly for none.

The paired writer restored 524 of 2,947 phrases, against 133 for mispaired supervision (difference 13.27 points, interval 11.84--14.70; Fig.~\ref{fig:visual-inference-writer}). Of these successes, 479 came from 880 targets visible intact in the premise. The writer recovered only 31 of 1,858 spans requiring at least one glyph absent from both pictured inputs. Neighboring-glyph counts alone recovered 196 in that subset. A selector calibrated on the internal validation cohort chose between the two complete image proposals using only their visible evidence. It recovered 651 official spans overall, compared with 524 for the neural proposal and 394 for the counts. Its gain over the neural writer was 4.31 points (interval 3.29--5.34). The selector recovered 170 absent-ink spans, below the count model's 196. On the newly written examples in panel c, changing the passage changed the missing drink correctly. Feeding the resulting premise and claim PNGs back into the writer reproduced each answer byte for byte.

Local Qwen3-8B Q8, given the original text and prescribed response format without training examples, recovered 73 of 105 relation sentences on the 35 complete premise triples. The paired visual writer recovered 59 (visual minus Qwen, $-13.3$ points, interval $-25.7$ to $-1.0$). On the same cases the calibrated writer recovered 26 missing phrases and Qwen recovered 18 (7.6 points, interval $-1.9$ to 18.1). The models differed in modality and training. The paired visual fit took 174.2 seconds after initialization from prior studies. This adaptation retained only five of 96 independent Chinese dictionary answers and gave zero exact natural DuReader answers out of 571. Short evidence-grounded restorations leave a substantial gap to a general language writer.
